\section{Summen- und Produktnotation}
\label{sec:fundamentals_summation_and_product_notation}

Bevor wir uns mit der Summen- und Produktnotation beschäftigen, wollen wir
zunächst den Begriff der \emph{Folge} einführen. Er ist für das Verständnis
der Summen- und Produktnotation von zentraler Bedeutung.

\begin{DefinitionBox}{Folge}{fundamentals_sequence}
  Eine \emph{Folge} \EnglishTerm{sequence} ist eine unendliche geordnete
  Anordnung von Zahlen oder Termen.
  \[
    (a_k)_{k\in\N}=(a_1,a_2,a_3,\ldots)
  \]
  Dabei bezeichnet $k$ den \emph{Index} und $a_k$ das dazugehörige
  Folgenglied. Die Punkte zeigen an, dass die Folge unbegrenzt fortgesetzt wird.
\end{DefinitionBox}

\begin{ExampleBox}{Folgen}{fundamentals_sequence_examples}
  \begin{enumerate}[label=(\alph*),leftmargin=2em,itemsep=0.7em]
    \item $a_k=k$ für $k\in\N$:
      \begin{ClozeArea}[30]{2}
        \[
          (a_k)_{k\in\N}=(1,2,3,4,\ldots)
        \]
      \end{ClozeArea}
    \item $b_k=2^{k-1}$ für $k\in\N$:
      \begin{ClozeArea}[30]{2}
        \[
          (b_k)_{k\in\N}=(1,2,4,8,16,\ldots)
        \]
      \end{ClozeArea}
    \item $c_k=(-1)^k$ für $k\in\N$:
      \begin{ClozeArea}[30]{2}
        \[
          (c_k)_{k\in\N}=(-1,1,-1,1,-1,\ldots)
        \]
      \end{ClozeArea}
  \end{enumerate}
\end{ExampleBox}

Auf dieser Folgenschreibweise bauen Summen- und Produktzeichen auf: Sie fassen
die Glieder einer Folge kompakt zu einer Summe beziehungsweise einem Produkt
zusammen.

\subsection{Summenzeichen}
\label{subsec:fundamentals_summation_notation}

\begin{DefinitionBox}{Summenschreibweise}{fundamentals_summation_notation}
  Die \emph{Summenschreibweise} \EnglishTerm{summation notation}
  \[
    \sum_{k=m}^{n}a_k
      =a_m+a_{m+1}+\dots+a_n
  \]
  fasst die Folgenglieder vom Index $m$ bis zum Index $n$ zu einer Summe
  zusammen. Dabei ist $k$ der \emph{Laufindex}, $m$ die untere und $n$ die
  obere Summationsgrenze.
\end{DefinitionBox}

\begin{ExampleBox}{Summen auswerten}{fundamentals_evaluating_sums}
  \begin{enumerate}[label=(\alph*),leftmargin=2em,itemsep=0.8em]
    \item\leavevmode
      \[
        \sum_{k=1}^{3}\frac1k
          =\Cloze{12cm}{1+\frac12+\frac13=\frac{11}{6}}
      \]
    \item\leavevmode
      \[
        \sum_{k=0}^{2}(k+1)^2
          =\Cloze{12cm}{(0+1)^2+(1+1)^2+(2+1)^2=14}
      \]
    \item\leavevmode
      \[
        \sum_{k=-2}^{3}k
          =\Cloze{12cm}{(-2)+(-1)+0+1+2+3=3}
      \]
  \end{enumerate}
\end{ExampleBox}



\subsubsection{Geschlossene Formen}
\label{subsubsec:fundamentals_closed_forms}

Eine \emph{geschlossene Form} gibt den Wert einer Summe durch einen endlichen
Term ohne Summenzeichen an. Eine solche Darstellung existiert nicht für jede
Summe. Eine \emph{Summenformel} ist eine Identität für eine bestimmte Art von
Summen. Kommt ihre rechte Seite ohne Summenzeichen aus, liefert sie eine
geschlossene Form. Drei wichtige geschlossene Formen sind:

\begin{enumerate}[label=\textbf{\arabic*.},leftmargin=2.3em,itemsep=1.2em]
  \item \textbf{Konstante Summe.}
    Für $n\in\N$ gilt
    \[
      \sum_{k=1}^{n}1
        =\Cloze{12cm}{\underbrace{1+1+\dots+1}_{n\text{-mal}}=n}.
    \]

  \item
    \begin{minipage}[t]{\linewidth}
      \textbf{Gaußsche Summenformel.}
      Für $n\in\N$ bezeichnen wir die Summe bis zum Index $n$ mit
      \[
        S_n:=\sum_{k=1}^{n}k=\Cloze{8cm}{1+2+\dots+n}.
      \]
      Schreibt man dieselbe Summe einmal vorwärts und einmal rückwärts, werden
      jeweils der erste und letzte, der zweite und vorletzte Summand usw.
      addiert:
      \begin{ClozeArea}[33]{12}
        \centering
        \begin{tikzpicture}[
  x=1cm,
  y=1cm,
  every node/.style={font=\normalsize},
  result/.style={text=THMBlue}
]
  \node[anchor=east] at (-6.1,1.25) {$S_n=$};
  \node at (-5.35,1.25) {$1$};
  \node at (-3.55,1.25) {$2$};
  \node at (-1.75,1.25) {$3$};
  \node at (0,1.25) {$\dots$};
  \node at (1.75,1.25) {$n-2$};
  \node at (3.55,1.25) {$n-1$};
  \node at (5.35,1.25) {$n$};

  \node at (-4.45,1.25) {$+$};
  \node at (-2.65,1.25) {$+$};
  \node at (-0.88,1.25) {$+$};
  \node at (0.88,1.25) {$+$};
  \node at (2.65,1.25) {$+$};
  \node at (4.45,1.25) {$+$};

  \node at (-7.45,0.15) {$+$};
  \node[anchor=east] at (-6.1,0.15) {$S_n=$};
  \node at (-5.35,0.15) {$n$};
  \node at (-3.55,0.15) {$n-1$};
  \node at (-1.75,0.15) {$n-2$};
  \node at (0,0.15) {$\dots$};
  \node at (1.75,0.15) {$3$};
  \node at (3.55,0.15) {$2$};
  \node at (5.35,0.15) {$1$};

  \node at (-4.45,0.15) {$+$};
  \node at (-2.65,0.15) {$+$};
  \node at (-0.88,0.15) {$+$};
  \node at (0.88,0.15) {$+$};
  \node at (2.65,0.15) {$+$};
  \node at (4.45,0.15) {$+$};

  \draw[line width=0.6pt] (-7.7,-0.45) -- (6.0,-0.45);

  \node[anchor=east] at (-6.1,-1.05) {$2S_n=$};
  \node[result] at (-5.35,-1.05) {$n+1$};
  \node[result] at (-3.55,-1.05) {$n+1$};
  \node[result] at (-1.75,-1.05) {$n+1$};
  \node[result] at (0,-1.05) {$\dots$};
  \node[result] at (1.75,-1.05) {$n+1$};
  \node[result] at (3.55,-1.05) {$n+1$};
  \node[result] at (5.35,-1.05) {$n+1$};

  \node at (-4.45,-1.05) {$+$};
  \node at (-2.65,-1.05) {$+$};
  \node at (-0.88,-1.05) {$+$};
  \node at (0.88,-1.05) {$+$};
  \node at (2.65,-1.05) {$+$};
  \node at (4.45,-1.05) {$+$};

\end{tikzpicture}

      \end{ClozeArea}
    \end{minipage}

    Jedes dieser Summandenpaare ergibt $n+1$. Damit ergibt sich
    \begin{ClozeArea}[33]{6}
      \[
        2S_n = n(n+1)
        \qquad \Longleftrightarrow \qquad
        S_n = \frac{n(n+1)}{2}.
      \]
    \end{ClozeArea}
    

    \begin{TheoremBox}{Gaußsche Summenformel}{fundamentals_gaussian_sum_formula}
      Für $n\in\N$ gilt
      \[
        \sum_{k=1}^{n}k=1+2+\dots+n=\frac{n(n+1)}{2}.
      \]
    \end{TheoremBox}

  \item
    \begin{minipage}[t]{\linewidth}
      \textbf{Endliche geometrische Summe.}
      Für $n\in\N_0$ bezeichnen wir die Summe der Potenzen von $q$ bis zum
      Index $n$ mit
      \[
        S_n:=\sum_{k=0}^{n}q^k=\Cloze{8cm}{1+q+q^2+\dots+q^n}.
      \]
    \end{minipage}

    \begin{minipage}{\linewidth}
      Bei der Herleitung sind zwei Fälle zu unterscheiden. \\

      \textbf{Fall 1:\enspace $q\neq1$.}\quad
      Durch Multiplikation mit $q$ und Subtraktion heben sich die mittleren
      Summanden auf:
      \begin{ClozeArea}[33]{20}
        \[
          \begin{array}{c@{\qquad}l@{\quad}l@{\quad}l}
            \phantom{\Longleftrightarrow}
              & \displaystyle S_n
              & \displaystyle =1+q+q^2+\dots+q^n,
              & \\[0.8em]
            \Longleftrightarrow
              & \displaystyle qS_n
              & \displaystyle =\phantom{1+{}}q+q^2+\dots+q^n+q^{n+1},
              & \\[0.8em]
            \Longleftrightarrow
              & \displaystyle S_n-qS_n
              & \displaystyle =1-q^{n+1},
              & \\[0.8em]
            \Longleftrightarrow
              & \displaystyle (1-q)S_n
              & \displaystyle =1-q^{n+1}
              & \displaystyle \Big|\,:(1-q), \\[0.8em]
            \Longleftrightarrow
              & \displaystyle S_n
              & \displaystyle =\frac{1-q^{n+1}}{1-q}.
              &
          \end{array}
        \]
      \end{ClozeArea}
      Da $q\neq1$ gilt, ist $1-q\neq0$; die Division durch $1-q$ ist daher
      zulässig.
    \end{minipage}

    \medskip
    \begin{minipage}{\linewidth}
      \textbf{Fall 2:\enspace $q=1$.}\quad
      Dann besteht die Summe aus $n+1$ Einsen:
      \begin{ClozeArea}[33]{4}
        \[
          S_n=\sum_{k=0}^{n}1^k
            =\underbrace{1+1+\dots+1}_{n+1\text{ Summanden}}
            =n+1.
        \]
      \end{ClozeArea}
    \end{minipage}

    \begin{TheoremBox}{Endliche geometrische Summenformel}{fundamentals_finite_geometric_sum_formula}
      Für $q\in\R$ und $n\in\N_0$ gilt
      \[
        \sum_{k=0}^{n}q^k
        =
        \begin{cases}
          \dfrac{1-q^{n+1}}{1-q}, & q\neq1,\\[0.8em]
          n+1, & q=1.
        \end{cases}
      \]
    \end{TheoremBox}
\end{enumerate}

\subsection{Produktzeichen}
\label{subsec:fundamentals_product_notation}

Analog zum Summenzeichen fasst das Produktzeichen zu multiplizierende
Folgenglieder kompakt zusammen.

\begin{DefinitionBox}{Produktschreibweise}{fundamentals_product_notation}
  Die \emph{Produktschreibweise} \EnglishTerm{product notation}
  \[
    \prod_{k=m}^{n}a_k
      =a_m\cdot a_{m+1}\cdots a_n
  \]
  fasst die Glieder $a_m$ bis $a_n$ zu einem Produkt zusammen. Wie beim
  Summenzeichen ist $k$ der Laufindex; $m$ und $n$ sind die untere und obere
  Grenze.
\end{DefinitionBox}

\begin{ExampleBox}{Produkte auswerten}{fundamentals_evaluating_products}
  \begin{enumerate}[label=(\alph*),leftmargin=2em,itemsep=0.8em]
    \item\leavevmode
      \[
        \prod_{k=1}^{5}k
          =\Cloze{12cm}{1\cdot2\cdot3\cdot4\cdot5=120}
      \]
    \item\leavevmode
      \[
        \prod_{k=1}^{4}k^2
          =\Cloze{12cm}{1^2\cdot2^2\cdot3^2\cdot4^2
            =1\cdot4\cdot9\cdot16=576}
      \]
    \item\leavevmode
      \[
        \prod_{k=1}^{3}\frac1k
          =\Cloze{12cm}{1\cdot\frac12\cdot\frac13=\frac16}
      \]
  \end{enumerate}
\end{ExampleBox}

